% ============================================================
%  Kapitel 5 — Diskussion und Fazit
% ============================================================
\section{Diskussion und Fazit}
\label{sec:diskussion}

\subsection{Kritisches Feld $B_c(0)$ und Dünnfilmeffekt}

Der gemessene Wert $B_c(0) = (2779 \pm 165)~\mathrm{mT}$ liegt um mehr als eine Größenordnung über dem Bulk-Literaturwert für Blei ($B_c^\mathrm{Bulk}(0) = 80~\mathrm{mT}$~\cite{Buckel2013}).
Dies ist eine direkte Konsequenz der Dünnen-Film-Geometrie: Bei einem Film der Dicke $d = 100~\mathrm{nm}$ in senkrechtem Magnetfeld ist die relevante Eindringtiefe die Pearl-Länge~\cite{Pearl1964}
\begin{equation}
    \Lambda = \frac{2\lambda_L^2}{d} = \frac{2 \times (39~\mathrm{nm})^2}{100~\mathrm{nm}} \approx 30~\mathrm{nm}.
\end{equation}
Da $\Lambda \ll d$, ist der Film im Regime $d \gg \Lambda$; tatsächlich ist jedoch die entscheidende Längenskala die Kohärenzlänge $\xi_\mathrm{Pb}(0) \approx 83~\mathrm{nm}$~\cite{Tinkham2004}.
Für $d < \xi$ gelten besondere geometrische Bedingungen, die das effektive kritische Feld stark erhöhen können.
Der beobachtete Wert ist mit der Literatur zu dünnen Pb-Filmen qualitativ konsistent~\cite{Tinkham2004}.

\subsection{Sprungtemperatur $T_c$}

Der Fit liefert $T_c = (6{,}861 \pm 0{,}039)~\mathrm{K}$, also etwa $0{,}34~\mathrm{K}$ unterhalb des Bulk-Wertes $T_c^\mathrm{Bulk} = 7{,}2~\mathrm{K}$.
Eine Reduktion von $T_c$ in dünnen Filmen ist bekannt und kann durch Größeneffekte sowie Verunreinigungen beim Aufdampfen verursacht werden~\cite{Tinkham2004}.
Die direkte Beobachtung, dass Messungen bei $T \geq 6{,}796~\mathrm{K}$ und $B = 0$ keine Supraleitung mehr zeigen, ist konsistent mit dem Fit-Ergebnis.

\paragraph{Anomalie bei $T = 6{,}901~\mathrm{K}$.}
Auffällig ist, dass bei der unmittelbar davor liegenden Messung ($T = 6{,}900~\mathrm{K}$) der Film bereits vollständig normalleitend war, während bei $T = 6{,}901~\mathrm{K}$ ein Übergang bei $B_c \approx 68~\mathrm{mT}$ gemessen wurde.
Da beide Temperaturen innerhalb der Messgenauigkeit des Lakeshore~340 ($\pm 5~\mathrm{mK}$) identisch sind, ist davon auszugehen, dass die Temperatur während des Messzyklus leicht driftete.
Dieses Datenpunkt liegt im Fit nahe dem Rand des definierten Bereichs und beeinflusst die Ergebnisse nur geringfügig.

\subsection{Kritischer Strom und GL-Exponent}

Der Exponent $n = 1{,}75 \pm 0{,}16$ aus dem $I_c(T)$-Fit ist mit dem theoretischen Ginzburg-Landau-Wert $n = 3/2$ für den Grenzfall $T \to T_c$ verträglich.
Der Temperaturbereich der $I_c$-Messungen ($6{,}5~\mathrm{K}$ bis $6{,}8~\mathrm{K}$) liegt bereits sehr nah an $T_c = 6{,}861~\mathrm{K}$, sodass das Potenzgesetz~\eqref{eq:ic_fit} gut anwendbar ist.
Der absolute Wert $I_c(0)$ aus dem Fit ($\approx 548~\mathrm{A}$, extrapoliert auf $T = 0$) ist eine sehr starke Extrapolation über einen Temperaturbereich, in dem keine Daten vorliegen, und daher wenig aussagekräftig.

\subsection{Limitierungen und nicht durchgeführte Versuche}

Mehrere Versuchsteile konnten nicht abgeschlossen werden:
\begin{itemize}
    \item \emph{Pickup-Spule (Meißner-Effekt, Suszeptibilität)}: Der Versuchsaufbau war defekt.
    \item \emph{YBCO $R(T)$ in $ab$- und $c$-Richtung}: Kein Stickstoff-Kryostat verfügbar.
    \item \emph{Quantitative Levitation}: Nur qualitativ beobachtet, keine Messung.
\end{itemize}
Die fehlenden Teile schränken die Möglichkeit ein, Typ-I- und Typ-II-Supraleitung direkt am gleichen Versuchstag zu vergleichen.

\subsection{Fazit}

Trotz der eingeschränkten Versuchsdurchführung konnten die zentralen Eigenschaften der Typ-I-Supraleitung am Bleifilm quantitativ charakterisiert werden.
Das parabolische Gesetz $B_c(T) = B_c(0)[1-(T/T_c)^2]$ wurde über fast den gesamten Temperaturbereich unterhalb $T_c$ bestätigt.
Der erhöhte Wert von $B_c(0) \approx 2779~\mathrm{mT}$ gegenüber dem Bulk ist eine klare Manifestation des Dünnfilmeffekts im senkrechten Magnetfeld.
Die Sprungtemperatur $T_c = 6{,}861~\mathrm{K}$ liegt leicht unterhalb des Bulk-Wertes, was für aufgedampfte Pb-Filme typisch ist.
Die $I_c(T)$-Messungen bestätigen qualitativ den erwarteten Potenzgesetzverlauf mit einem GL-konsistenten Exponenten.
Das Levitationsexperiment demonstrierte anschaulich den Meißner-Effekt und das Flux-Pinning in YBCO.
